\chapter{El poder de la incitación}\label{sec:el-poder-de-la-incitacion}

Era bastante fácil poner a una multitud en un estado de frenesí, siempre y cuando supieras lo que querían oír. La gente que acudió a mi manifestación hoy era un grupo temeroso. Las imágenes de la batalla revelaron que los terranos eran unos charlatanes, y las masas no olvidarían su engaño de siglos. Se habían enconado dentro de nuestras filas, conspirando, manipulando nuestra política con fines desconocidos, propagando su vil cultura. Una especie tan guerrera como la humanidad era como un tumor que consumía todo lo que era natural y bueno, a menos que fuera extirpado.

Tenía fe en que la gente reconocería al monstruo que acechaba entre ellos. Mi discurso de hoy solo estaba dando a la ciudadanía el empujón que necesitaba. Si los plebeyos dirigían su ira hacia los humanos, eso los obligaría a actuar. Los soldados terranos aniquilando a civiles en la calle serían bien recibidos por los medios de comunicación y, combinados con las noticias de la masacre de los devoradores que llegarían en cualquier momento, la protesta podría ser suficiente para expulsarlos de la Federación.

—Cuando era niña, la humanidad era una especie que se destacaba en nuestra clase de educación cívica como la de los mejores diplomáticos. Los guardianes de la paz —hice una pausa y mi mirada recorrió a la multitud—. Pero hace poco nos dimos cuenta de que todo eso era mentira. Son brutales, sedientos de sangre; no son como nosotros y no comparten nuestros valores. Los valores que mantienen unida a nuestra Federación.

Se escucharon algunos aplausos entre los presentes, pero la mayoría de los espectadores parecían ansiosos. Ese era exactamente el tipo de sentimiento que esperaba despertar; no hay motivación más poderosa que el miedo. Sabía que tenía toda su atención y que estarían pendientes de cada una de mi palabras.

—Amigos míos, tengan la seguridad de que comparto su sorpresa y confusión ante la noticia. Pero de una cosa estoy segura: debemos actuar ahora. Los humanos merodean por nuestras calles y es cuestión de tiempo que cedan a sus instintos. ¿Cuántos de ustedes le han dado la bienvenida a uno en sus hogares? ¿Han enviado a sus hijos a la escuela con uno?

Expresiones de horror, murmullos agitados; esa era la respuesta que esperaba. La mayoría de la gente allí había tenido contacto con un terrano al menos una vez. Sin duda, algunos oyentes dudarían en juzgar a la humanidad, pero un recordatorio del peligro mortal que representaban los humanos debería asustarlos y hacerlos entrar en razón.

 —¿Qué vamos a hacer? ¿Qué van a hacer ustedes? —mi voz se convirtió en un gruñido bajo—. Díganlo conmigo, lo suficientemente fuerte para que se escuche hasta en la Tierra: ¡Los humanos no son bienvenidos aquí!

—¡Los humanos no son bienvenidos aquí! —gritó la multitud.

Una sonrisa se dibujó en mi rostro. —Es un sonido hermoso. Ahora, la embajada terrana está en esta misma calle, a solo unos minutos de aquí. Contamina nuestra capital con su presencia. ¿Por qué no hacen oír sus voces allá? ¡Vamos a recuperar lo que es nuestro!

En respuesta se escucharon vítores de aprobación y observé con satisfacción cómo la gente volvía la vista hacia el complejo cerrado al final de la calle. No fue casualidad que hubiera elegido un lugar al aire libre para este evento, a solo unos minutos a pie de la embajada terrestre. Los humanos no tendrían tiempo de prepararse. Una multitud furiosa de miles de personas superaría rápidamente su seguridad y entonces se verían obligados a tomar medidas más drásticas.

Saqué mi holopad del bolsillo mientras la multitud se alejaba. El siguiente punto de mi agenda era ponerme en contacto con el general Kilon y averiguar qué había pasado con la misión de rescate de los terranos. Gracias a mi intromisión, no tendrían más opción que recurrir a la violencia. El general a menudo simpatizaba con los humanos, probablemente por gratitud por haberle salvado la vida. Si algo podía hacerle cambiar de opinión, sería la masacre de un planeta entero. Ya no vería a la humanidad como salvadores, sino como los monstruos que realmente eran.

Mi llamada a la flota no fue recibida o no fue respondida. Una pequeña semilla de duda se plantó en mi mente mientras consideraba la posibilidad de que los terranos los hubieran aniquilado en represalia. No había forma de que supieran lo que había hecho una vez que las naves se perdieron, por supuesto. Pero ¿y si habían culpado a la flota de todos modos o simplemente no querían testigos de su masacre?

La idea me heló los huesos. Tenía la esperanza de que el general estuviera ocupado con otros asuntos y me devolviera la llamada pronto. Lo que fuera que les sucediera a las naves humanas no me importaba, pero la pérdida de nuestros soldados sería trágica. Tal vez mis acciones habían sido un poco descuidadas.

Pero ya era demasiado tarde para dar marcha atrás. Era inevitable que algunas personas salieran lastimadas al tratar con una especie tan agresiva, pero su sacrificio era necesario por el bien común. Los acontecimientos de hoy eran el ejemplo perfecto de ese principio. Dudaba que el caos que se estaba desatando en la embajada fuera incruento.

Uno de mis colaboradores había avisado a Federation News Central, diciéndoles que mantuvieran un equipo de filmación apostado en la embajada terrana. De lo contrario, los medios podrían haberse perdido el comienzo de la manifestación. Cambié a la transmisión en vivo en mi holopad y me alegré de ver que la cobertura de la protesta ya estaba en marcha.

\textit{Dos centinelas humanos se paseaban por el interior de las puertas y gritaban a todo el mundo que se hiciera a un lado. La multitud los abucheaba como respuesta, lanzando piedras y otros proyectiles por encima de las barreras.}

\textit{Un joven reportero xanik señaló la escena que se desarrollaba detrás de él: —Como pueden ver, la situación en la embajada terrestre se está agravando. La presidenta Ula pronunció un breve discurso esta tarde, cargado de retórica antihumana. Sus partidarios se sintieron inspirados a actuar y, alentados por la presidenta, se dirigieron al complejo. Parece que la confrontación es inminente.}

\textit{Algunos manifestantes comenzaron a escalar los muros y los guardias terrestres les apuntaron con sus armas en respuesta. Sin inmutarse, los escaladores se lanzaron al otro lado.}

Me reí para mis adentros. En el momento en que los humanos abrieran fuego contra civiles desarmados, todos los verían como yo. Este incidente sería transmitido en vivo a toda la galaxia y se reproduciría una y otra vez por días, tal como lo fue la nanobomba de nanitos.

\textit{Más civiles llegaron al lugar. Los guardias retrocedieron unos pasos, con las armas aún en alto. Sus dedos se cernían sobre el gatillo, listos para disparar en cuanto un manifestante avanzara. La multitud cerró filas y luego cargó en conjunto.}

\textit{En lugar del sonido de los disparos, lo único que oí fue una voz humana que gritaba que se retiraran. Inexplicablemente, los guardias enfundaron sus armas y se retiraron al interior del edificio.}

—¡¿Por qué no disparan?! —grité—. ¡Esto no está bien! ¡Se supone que los humanos son asesinos!

Mis colaboradores, que estaban apiñados cerca, se sorprendieron por mi arrebato. Me miraban como si hubiera perdido la cabeza.

Un joven llamado Radi parecía particularmente preocupado. —¿Quieres que disparen? ¿Estás bien?

Lo miré con enojo. —¡Ocúpate de tus asuntos! Estoy perfectamente bien.

Los manifestantes habían entrado en la embajada sin que yo estuviera prestando atención y ya no eran visibles para la cámara. El presentador dijo algo sobre una situación de rehenes, pero no lo registré. ¿Cómo pudo haber sucedido esto? ¿Una especie agresiva y militarizada rindiéndose sin luchar? ¡Aquello fue un desastre absoluto!

Con sus diplomáticos capturados, los humanos serían vistos como las víctimas de la historia. Podía oírlo ahora, un terrano en las noticias diciéndoles a los espectadores que ellos no eran los violentos. Eso distraería la atención de la atrocidad de sus armas y, en cambio, la controversia se centraría en mí, quien incitó el motín.

Este era el tipo de escándalo que podría paralizar la carrera de un político promedio. Pero yo no era un representante menor; sin duda, una presidenta tan popular como yo podría capear el temporal. En lugar de disculparme por mis acciones, redoblaría mi apuesta. Con el alma misma de la Federación en juego, renunciar a mi misión de desenmascarar a la humanidad no era una opción.

Con suerte, mis esfuerzos por sabotear su misión de rescate habían tenido más éxito. Un solo incidente era todo lo que necesitaba para demostrar mi punto, y sabía que los humanos estaban destinados a cometer errores en algún momento. Estaba en su naturaleza, después de todo.
